\section{Bottom-Up trace}
\label{app:trace}
Table~\ref{tab:trace} traces the algorithm on \texttt{baskets.txt}, a small database with
the four transactions $T_1 = \{1,2,3\}$, $T_2 = \{1,2\}$, $T_3 = \{1,3\}$, and
$T_4 = \{2,3\}$, at $\sigma_{\min} = 2$. The search starts from the tidsets
$t(1) = \{1,2,3\}$, $t(2) = \{1,2,4\}$, and $t(3) = \{1,3,4\}$. It records exactly six
frequent itemsets and prunes $\{1,2,3\}$, which only $T_1$ contains.

\begin{table}[t]
\centering
\caption{Bottom-Up trace on \texttt{baskets.txt} with $\sigma_{\min} = 2$}
\label{tab:trace}
\begin{tabular}{llll}
\hline
Prefix & Intersection & Tidset & Result \\
\hline
$\emptyset$ & -- & $t(1) = \{1,2,3\}$ & record $\{1\}$, $\sigma = 3$ \\
$\{1\}$ & $t(1) \cap t(2)$ & $\{1,2\}$ & keep $2$ \\
$\{1\}$ & $t(1) \cap t(3)$ & $\{1,3\}$ & keep $3$ \\
$\{1\}$ & -- & $t(12) = \{1,2\}$ & record $\{1,2\}$, $\sigma = 2$ \\
$\{1,2\}$ & $t(12) \cap t(13)$ & $\{1\}$ & prune, $1 < 2$ \\
$\{1\}$ & -- & $t(13) = \{1,3\}$ & record $\{1,3\}$, $\sigma = 2$ \\
$\emptyset$ & -- & $t(2) = \{1,2,4\}$ & record $\{2\}$, $\sigma = 3$ \\
$\{2\}$ & $t(2) \cap t(3)$ & $\{1,4\}$ & keep $3$ \\
$\{2\}$ & -- & $t(23) = \{1,4\}$ & record $\{2,3\}$, $\sigma = 2$ \\
$\emptyset$ & -- & $t(3) = \{1,3,4\}$ & record $\{3\}$, $\sigma = 3$ \\
\hline
\end{tabular}
\end{table}

\section{Test groups}
\label{app:tests}
Both variants go through the following groups of tests.

\begin{itemize}
\item Hand-derived results: eight small databases, including an empty database, a
database with no frequent itemset, duplicate transactions, thresholds just above and
just below an itemset's support, and an image database with empty transactions.
\item Comparison with SPMF's Eclat~\cite{spmf2016}: eight runs on seven small databases
and one run on T25I10D10K (9\,976 transactions, 929 items), where SPMF reports 14\,782
itemsets of 1 to 13 items at $\mathit{minsup} = 0.008$. Both variants match 100\% on all
nine runs, with no missing, extra, or wrong-support itemset.
\item Comparison with a brute-force counter that enumerates every subset, on 10 random
databases generated with seed 42 for each variant. The optimized variant is also tested
on eight databases with 0, 1, 63, 64, 65, 127, 128, and 129 transactions, sizes that
straddle the boundary of a 64-bit bitset word. Half of them contain groups of items that
always occur together, so that perfect extensions appear. Each database is mined under 18
configurations that turn the techniques of Section~\ref{sec:optimized} on and off.
\item Type stability: \texttt{Test.@inferred} on \texttt{mine}, \texttt{mine\_packed},
and the search functions of both variants.
\item Reading, writing, and the command line: the byte-level reader is checked against a
line-based parser on edge-case inputs (\texttt{Int} overflow, unexpected characters,
blank lines, Windows line endings) and on every small database. The writer's output is
compared byte for byte with a reference formatter, and the tests also cover argument
errors and process exit codes.
\end{itemize}

\section{Supplementary measurements}
\label{app:measurements}

\subsection{Correctness on small databases}
\label{app:toy}
Table~\ref{tab:toy-spmf} shows the results on seven small databases, with $N$ and
$\sigma_{\min}$ as defined in Section~\ref{sec:protocol}. The ``Match'' column gives the
number of itemsets whose support agrees with SPMF out of the total number of distinct
itemsets. The baseline reaches 100\% in every comparison, with no missing, extra,
duplicate, or wrong-support itemset. On \texttt{disjoint}, both programs return an empty
map. The \texttt{boundary} database is mined at two thresholds to check itemsets that lie
just above and just below the threshold, and each threshold is chosen so that SPMF and
exact decimal arithmetic agree on $\sigma_{\min}$. The test suite checks every Eclat
variant against the stored SPMF outputs each time it runs.

\begin{table}[t]
\centering
\caption{Full itemset and support comparison of the baseline Eclat with SPMF on small
databases}
\label{tab:toy-spmf}
\begin{tabular}{lrrr}
\hline
Data & $N$ & $\sigma_{\min}$ & Match \\
\hline
Baskets & 4 & 2 & 6/6 \\
Duplicates & 3 & 2 & 3/3 \\
Disjoint & 3 & 2 & 0/0 \\
Dense & 2 & 2 & 7/7 \\
Groceries & 5 & 2 & 5/5 \\
Boundary & 4 & 2 & 2/2 \\
Boundary & 4 & 1 & 4/4 \\
Nested & 4 & 2 & 3/3 \\
\hline
\end{tabular}
\end{table}

\subsection{Measurement noise}
\label{app:noise}
Relative to the median, the interquartile range of the runtime campaign is at most 2.9\%
for the baseline and 1.9\% for the optimized variant. It is wider for the Java programs,
reaching 11.3\% for SPMF and 13.9\% for PrePost+, although only one SPMF point and three
PrePost+ points exceed 10\%. To estimate drift between sessions, we compared this session
with three others on the same machine. PrePost+, which runs the same Java command every
time, differs by $+3.0\%$, $+1.9\%$, and $-2.3\%$ in geometric mean, and by at most
14.4\% at a single point. The optimized variant, measured in another session at the same
revision, differs by $-1.5\%$ on average and by at most 5.6\%. Because the four programs
in the tables ran in the same session, a ratio between two of them carries only
within-session noise. We treat a ratio as a real difference when it exceeds the sum of
the two largest interquartile ranges, which comes to about 5\% between the two Julia
variants and about 16\% when one side is a Java program. Drift between sessions only
matters when comparing with another campaign.

At the highest minsup of Chess, Mushroom, Accidents, and Retail, the baseline is only
1.09, 1.06, 1.27, and 1.46 times slower than the optimized variant. All four ratios
exceed the 5\% margin, so the differences are real, but they are small. On Chess and
Mushroom the gap is a few percent of a run that consists mostly of start-up and file
reading. On Accidents and Retail the absolute gap is only 0.21\,s and 0.13\,s, about as
long as start-up and file reading themselves (some 0.21\,s on Chess at minsup 0.9, a
point that does little else). These small ratios therefore do not mean that the two
algorithms are equally fast.

\subsection{Runtime ratios over the whole grid}
\label{app:ratios}
Table~\ref{tab:runtime-ratio} summarizes each benchmark's runtime curves of
Fig.~\ref{fig:runtime} with one number per program.

\begin{table}[t]
\centering
\small
\setlength{\tabcolsep}{4pt}
\caption{Geometric mean over every point of each benchmark of the \emph{process} time
ratio: the time of the program named in the column header divided by the time of the
optimized variant. ``SPMF'' is SPMF's Eclat. A value above 1 means slower than the
optimized variant.}
\label{tab:runtime-ratio}
\begin{tabular}{lrrr}
\hline
Data & Baseline & SPMF & PrePost+ \\
\hline
Chess & 2.83 & 10.33 & 2.08 \\
Mushroom & 2.73 & 6.01 & 2.22 \\
Retail & 5.89 & 4.34 & 5.01 \\
T10I4D100K & 13.14 & 4.53 & 4.76 \\
Accidents & 2.62 & 10.76 & 2.39 \\
\hline
\end{tabular}
\end{table}

\subsection{Intersection work of the baseline}
\label{app:join-work}
To measure the baseline's work, we reran its search on one thread and counted the
intersections (joins), how many of them produced a frequent itemset, and the combined
length of the two tidsets summed over all intersections, which we call merged elements
(Table~\ref{tab:join-work}). Taking the time at the lowest minsup, subtracting the time at
the highest minsup (mostly start-up and file reading), and dividing by the difference in
merged elements gives 2.0 to 6.9\,ns per element across the five benchmarks. This is the
same order of magnitude as the 9.0 to 10.0\,ns measured in the transaction-length
experiment of Section~\ref{sec:width}. The intersection counts come from a separate run,
outside the timed runs.

\begin{table}[t]
\centering
\small
\setlength{\tabcolsep}{3pt}
\caption{Work done by the baseline at the lowest minsup of each benchmark. ``Success'' is
the share of intersections that produce a frequent itemset, ``Elements'' is the combined
length of the two tidsets over all intersections (billions), and ns/el is the time per
merged element after subtracting the time at the highest minsup.}
\label{tab:join-work}
\begin{tabular}{lrrrr}
\hline
Data / minsup & Joins & Success & Elements & ns/el \\
\hline
Chess / 0.6 & 326\,545 & 78.1\% & 1.43 & 2.0 \\
Mushroom / 0.1 & 722\,721 & 79.5\% & 1.66 & 3.6 \\
Retail / 0.001 & 3\,932\,582 & 0.14\% & 2.17 & 5.7 \\
T10I4D100K / 0.00025 & 4\,192\,075 & 2.3\% & 1.48 & 6.9 \\
Accidents / 0.65 & 2\,333 & 46.1\% & 1.22 & 4.8 \\
\hline
\end{tabular}
\end{table}

On Retail 0.001, each frequent itemset costs about 286\,000 merged elements. On
T10I4D100K, even minsup 0.01 takes 70\,127 intersections and 297.7 million merged elements
to find just 10 two-item sets, so the baseline is already 8.5 times slower than the
optimized variant at the first point.

\subsection{A model of output growth}
\label{app:output-model}
Model~\eqref{eq:output-growth} is the logarithm of a power law,
$F \approx e^{c}\,\mathrm{minsup}^{-s}$, and a simple model of the data explains why a
power law fits. Suppose each of the $|I|$ items occurs in a transaction independently with
probability $d$, the density of Table~\ref{tab:datasets}. A $k$-itemset then occurs in a
fraction $d^k$ of the transactions, so it is frequent when $d^k \ge \mathrm{minsup}$, that
is, when $k \le k^* = \log(1/\mathrm{minsup}) / \log(1/d)$. The number of itemsets of
length up to $k^*$ grows like $|I|^{k^*}$, so
\begin{equation}
  \log F \approx k^* \log|I| = s \log(1/\mathrm{minsup}), \qquad
  s \approx \frac{\log|I|}{\log(1/d)}.
  \label{eq:output-slope}
\end{equation}
Estimate~\eqref{eq:output-slope} gives 1.5 for T10I4D100K and 1.3 for Retail, close to
the measured 1.6 and 1.7. It also predicts that denser data grows faster, since a larger
$d$ shrinks $\log(1/d)$. On the dense databases it gives only 2.3 to 6.1, because their
items are correlated and their frequent items occur far more often than the average
density $d$.

\subsection{Output size and runtime}
\label{app:count-runtime}
On Retail, most of the two million root intersections fail without producing any itemset
(Section~\ref{sec:optimized}). The optimized variant avoids them by counting pairs and
runs Retail 0.001 in 0.38\,s. Even so, it takes 0.41\,s on Mushroom 0.1 with
574\,431 itemsets but 0.96\,s on T10I4D100K 0.00025 with only 96\,959, because the
tidsets of T10I4D100K span 100\,000 transactions and most intersections at depth 2 fail.
The count dominates the runtime only when it is very large: on Mushroom 0.05, before the
radix sort was added, sorting and writing 3\,755\,511 itemsets took 89\% of the time and
mining only 11\% (Section~\ref{sec:optimized}). Conversely, when the relative minsup is
fixed, the count barely changes and the runtime grows with $N$, as
Section~\ref{sec:scalability} shows.

\subsection{Other factors in the scalability subsets}
\label{app:scal-factors}
Two factors besides $N$ could change the amount of work between subsets, and we measured
both. The first is the number of distinct items, which grows with $N$ from 343 to 468 on
Accidents and from 10\,208 to 16\,470 on Retail. The number of frequent items, however,
barely changes: exactly 18 in all five Accidents subsets and between 2\,100 and 2\,270 on
Retail. The new items are all rare and drop out in the first counting pass. The second
factor is the very low absolute threshold on Retail 10\%, where $\sigma_{\min} = 9$. In a
small sample, the supports of itemsets that sit just below the threshold fluctuate more,
and since there are very many such itemsets, more of them rise above the threshold than
fall below it. The 10\% subset therefore has the most frequent items of the five subsets
(2\,270) and 8\,895 frequent itemsets, 17\% more than the 7\,589 of the full database.
This extra work is small compared with the growth in $N$: the point farthest from the
fitted line is off by only 4\,ms for the optimized variant on Retail and by 0.12\,s for
the baseline, which is 1\% of its time on the full database.

\subsection{Search trees of the two width protocols}
\label{app:width-tree}
The two protocols of Section~\ref{sec:width} shape the search tree differently. With
minsup fixed, the tree gains a level at $w = 25$ (pairs) and again at $w = 40$ (triples),
and it widens as well: the number of intersections grows from 4\,950 to 166\,650 and then
to 4\,087\,798. With the output matched, the tree has the same size at every $w$, and only
the cost of each intersection changes.

\section{Tradeoffs and tests of the two proposals}
\label{app:proposals}
\emph{Proposal 1, tradeoff.} The CLI would gain a second reader and a second layout
builder alongside the general ones, and the two would have to stay equivalent. The
item-code cutoff also needs a fallback path because SPMF files allow large, sparse codes,
although the codes of all five benchmarks are small enough for an array.
\emph{Tests.} The new layout must equal \texttt{bitset\_class} on seeded random databases,
including $N = 63, 64, 65$, code 0, and codes above the cutoff. The CLI output must stay
byte-identical to the current output at all 33 points and match SPMF's full map. The
layout share, slope, and peak RSS would then be measured again, and the change kept only
if it gains at least 10\% at the targeted points without slowing any other point by more
than 10\%.

\emph{Proposal 2, tradeoff.} Files written in search order could no longer be compared
byte for byte across variants or across code revisions, which the package interface
currently guarantees through its fixed order. The comparator reads both files into maps,
so it would not be affected, but a benchmark would have to record which order it used,
because CLI time would then exclude the sort, as it already does for the SPMF programs.
Since the measured gain is concentrated at the dense points, and above all at Mushroom
0.05, the package keeps the fixed order. \emph{Tests.} A search-order file, once sorted
externally, must be byte-identical to the default output, and the full map must match
SPMF at every grid and stress point. Two runs of the same code must produce the same
bytes. A five-repeat campaign of both orders on Chess 0.6, Mushroom 0.1, and Mushroom 0.05
would test the prediction, which fails if $R$ on Mushroom 0.05 stays at 1.0 or above.
